\section{Related Works}
\label{sec:related}

\paragraph{General purpose features from Vision Transformers.} Transformers have been used extensively across multiple domains as general-purpose feature extractors \cite{brown2020language,radford2019language,chowdhery2022palm,touvron2023llama,devlin2019bert, raffel2023exploring}.
Vision Transformers \cite{dosovitskiy2021image} (ViTs) pre-trained via supervised learning \cite{vaswani2023attention, touvron2022deit,kirillov2023segment} or self-supervised learning \cite{zhou2021ibot,he2022masked,caron2021emerging,oquab2023dinov2} have demonstrated strong generalizability to various downstream visual tasks, even without fine-tuning. However, we show that ViTs trained with diverse training objectives exhibit commonly observed noise artifacts in their output feature maps. These artifacts are often overlooked in practice because their presence cannot be simply reflected by image \textit{classification} accuracy. Thus, our work focuses on evaluating pre-trained ViTs for \textit{dense recognition} tasks such as segmentation, depth estimation, and object discovery. We demonstrate how these artifacts adversely affect dense recognition tasks, thereby motivating our method to mitigate them.

\paragraph{ViT artifacts.} Our work studies the noise artifacts in ViTs, an issue that has been previously observed but often remains unexplored. These artifacts manifest as noisy attention maps in supervised ViTs (\ie, ViTs do not attend to objects of interest well) \cite{caron2021emerging,chen2021vision}. Concurrently with our study, two recent studies similarly have also identified artifacts in self-supervised ViTs \cite{yang2023emernerf,darcet2023vision}. Specifically, \cite{darcet2023vision} describe these as ``high-norm'' patches in low-informative background regions, hypothesizing their occurrence is limited to large (\eg ViT-\textit{large} or greater) and sufficiently trained ViTs. However, our analysis indicates that this may not be the \textit{full picture}, as we observe similar artifacts in small or base ViTs that cannot be easily identified by extremely high feature norm values. Instead, we find a strong correlation between the presence of artifacts and the use of positional embeddings in ViTs. This finding suggests that artifacts are not strictly confined to certain model sizes or training scales but are more fundamentally linked to the inherent design of ViTs.
Moreover, unlike the method proposed by \cite{darcet2023vision} that retrains ViTs with register tokens \cite{goyal2023think, xiao2023efficient} from scratch,
our approach directly denoises pre-trained models without retraining. Users can \textit{dynamically} enable or disable the plugged-in denoiser as needed. Lastly, we note that some \textit{weak} artifacts still exist in DINOv2 models trained with registers \cite{darcet2023vision} (see \cref{fig:teaser} DINOv2-reg and appendix), and our \ours can effectively denoise them, improving the performance of DINOv2-reg.

